\section{结论}

多年累积的卫星海洋遥感大数据对于自主水下航行器的意义，不仅在于提供一张指引大方向的先验底图，更在于能够生成一套经真实误差标定且独立于机载规划器搜索过程的留出想象海洋，使航行器拥有衡量自身计算价值的客观标尺而不被规划器的内部自评泡沫所误导。本文围绕东海赤潮应急巡测中走、深想与停下复测的时间与能量分配问题，将48个随现场采样局地化集合同化更新的想象海洋划分为32个搜索用海洋与16个留出海洋，分别提出了下水前在母船或岸基计算机上运行的出发前全任务预演方法以及航行途中逐周期运行的在线留出估值单步元推理方法。在2018至2026年东海3197天卫星数据构建的422个赤潮窗口、四种机载硬件及两种真值模型上的实验证实，留出海洋评判在横跨三个半数量级的推演预算下始终与真实航线价值保持在0.005以内，开启高功耗芯片深想仅在基础算力过慢时有益而停下复测始终得不偿失，且出发前预演与在线留出估值在无需任何历史巡测任务的条件下均取得了与依赖历史任务方法一致、并显著优于经典单步元推理和启发式巡测规则的表现。未来工作将把基于留出想象海洋的计算价值评判框架拓展至三维斜压海流预报集合与多台自主水下航行器协同赤潮追踪的海试研究中。
